\section{Yang--Mills 梯度流}
\label{sec:3}
\subsection{流方程与微扰展开}
Yang--Mills 梯度流是把 Yang--Mills 作用量泛函视为势的高度、由梯度流所生成的规范场位形
变形。明确地说，对规范势~$A_\mu(x)$，流动定义为~\cite{Luscher:2010iy,Luscher:2011bx}
\begin{equation}
   \partial_tB_\mu(t,x)=D_\nu G_{\nu\mu}(t,x)
   +\alpha_0D_\mu\partial_\nu B_\nu(t,x),\qquad
   B_\mu(t=0,x)=A_\mu(x),
\label{eq:(3.1)}
\end{equation}
其中 $t$~是流时间，$G_{\mu\nu}(t,x)$ 是流动场的场强，
\begin{equation}
   G_{\mu\nu}(t,x)
   =\partial_\mu B_\nu(t,x)-\partial_\nu B_\mu(t,x)
   +[B_\mu(t,x),B_\nu(t,x)],
\label{eq:(3.2)}
\end{equation}
而作用在规范场上的协变导数为
\begin{equation}
   D_\mu=\partial_\mu+[B_\mu,\cdot].
\label{eq:(3.3)}
\end{equation}
式~\eqref{eq:(3.1)} 右边的第一项是泛函空间中的“梯度”，即
$-g_0^2\delta S/\delta B_\mu(t,x)$，其中 $S$~是流动场的 Yang--Mills 作用量泛函。注意，
由于 $D_\nu G_{\nu\mu}(t,x)=\Delta B_\mu(t,x)+O(B^2)$，式~\eqref{eq:(3.1)}
是一类扩散方程，正的~$t$ 的流有效地压制位形中的高频模式。式~\eqref{eq:(3.1)} 右边含
参数~$\alpha_0$ 的第二项是使微扰展开良定义的“规范固定项”。然而可以证明，任何规范不变量
都与~$\alpha_0$ 无关。在下节实际的微扰计算中我们采用表达式最简单的“费曼规范”
$\alpha_0=1$。

式~\eqref{eq:(3.1)} 的形式解由~\cite{Luscher:2010iy,Luscher:2011bx}给出：
\begin{equation}
   B_\mu(t,x)
   =\int\mathrm{d}^Dy\left[
   K_t(x-y)_{\mu\nu}A_\nu(y)
   +\int_0^t\mathrm{d}s\,K_{t-s}(x-y)_{\mu\nu}R_\nu(s,y)
   \right],
\label{eq:(3.4)}
\end{equation}
其中\footnote{在本文中我们始终使用如下缩写：
\begin{equation}
   \int_p\equiv\int\frac{\mathrm{d}^Dp}{(2\pi)^D}.
\label{eq:(3.5)}
\end{equation}
}
\begin{equation}
   K_t(x)_{\mu\nu}=\int_p\frac{\mathrm{e}^{ipx}}{p^2}
   \left[(\delta_{\mu\nu}p^2-p_\mu p_\nu)\mathrm{e}^{-tp^2}
   +p_\mu p_\nu \mathrm{e}^{-\alpha_0tp^2}\right]
\label{eq:(3.6)}
\end{equation}
是热核，而
\begin{equation}
   R_\mu=2[B_\nu,\partial_\nu B_\mu]
   -[B_\nu,\partial_\mu B_\nu]
   +(\alpha_0-1)[B_\mu,\partial_\nu B_\nu]
   +[B_\nu,[B_\nu,B_\mu]]
\label{eq:(3.7)}
\end{equation}
表示非线性相互作用项。对方程~\eqref{eq:(3.4)} 迭代求解，便得到流动
场~$B_\mu(t,x)$ 以初值~$A_\nu(y)$ 表出的微扰展开。

对费米子场也可以考虑类似的流~\cite{Luscher:2013cpa}。就我们的目的而言，并不要求费米子场
的流是原始费米子作用量的梯度流。文献~\cite{Luscher:2013cpa}中引入的一个可行选择是
\begin{align}
   &\partial_t\chi(t,x)=\left[\Delta-\alpha_0\partial_\mu B_\mu(t,x)\right]
   \chi(t,x),\qquad
   \chi(t=0,x)=\psi(x),
\label{eq:(3.8)}\\
   &\partial_t\Bar{\chi}(t,x)
   =\Bar{\chi}(t,x)
   \left[\overleftarrow{\Delta}
   +\alpha_0\partial_\mu B_\mu(t,x)\right],
   \qquad\Bar{\chi}(t=0,x)=\Bar{\psi}(x),
\label{eq:(3.9)}
\end{align}
其中
\begin{align}
   &\Delta=D_\mu D_\mu,\qquad D_\mu=\partial_\mu+B_\mu,
\label{eq:(3.10)}\\
   &\overleftarrow{\Delta}=\overleftarrow{D}_\mu\overleftarrow{D}_\mu,
   \qquad\overleftarrow{D}_\mu\equiv\overleftarrow{\partial}_\mu-B_\mu.
\label{eq:(3.11)}
\end{align}
上述流方程的形式解为
\begin{align}
   \chi(t,x)&=\int\mathrm{d}^Dy\,
   \left[
   K_t(x-y)\psi(y)
   +\int_0^t\mathrm{d}s\,K_{t-s}(x-y)\Delta'\chi(s,y)
   \right],
\label{eq:(3.12)}\\
   \Bar{\chi}(t,x)&=\int\mathrm{d}^Dy\,
   \left[
   \Bar{\psi}(y)K_t(x-y)
   +\int_0^t\mathrm{d}s\,\Bar{\chi}(s,y)\overleftarrow{\Delta}'K_{t-s}(x-y)
   \right],
\label{eq:(3.13)}
\end{align}
其中
\begin{equation}
   K_t(x)\equiv\int_p \mathrm{e}^{ipx}\mathrm{e}^{-tp^2}=\frac{\mathrm{e}^{-x^2/4t}}{(4\pi t)^{D/2}},
\label{eq:(3.14)}
\end{equation}
且
\begin{align}
   \Delta'&\equiv(1-\alpha_0)\partial_\mu B_\mu
   +2B_\mu\partial_\mu+B_\mu B_\mu,
\label{eq:(3.15)}\\
   \overleftarrow{\Delta}'&
   \equiv-(1-\alpha_0)\partial_\mu B_\mu
   -2\overleftarrow{\partial}_\mu B_\mu+B_\mu B_\mu.
\label{eq:(3.16)}
\end{align}

上述流的初值 $A_\mu(x)$、$\psi(x)$
和~$\Bar{\psi}(x)$ 是参与泛函积分的量子场。于是流动场的量子关联函数这样得到：把流动场用
原始的未流动场（初值）表出，然后对后者取泛函平均。
式~\eqref{eq:(3.4)}、\eqref{eq:(3.12)} 和~\eqref{eq:(3.13)} 提供了完成此事的显式方法。
例如，在最低阶（树图级）近似下有
\begin{equation}
   \left\langle B_\mu^a(t,x)B_\nu^b(s,y)\right\rangle
   =g_0^2\delta^{ab}\delta_{\mu\nu}\int_p\mathrm{e}^{ip(x-y)}\frac{\mathrm{e}^{-(t+s)p^2}}{p^2},
\label{eq:(3.17)}
\end{equation}
此即在 $\lambda_0=\alpha_0=1$ 的“费曼规范”，其中 $\lambda_0$
是常规的规范固定参数。类似地，树图级近似下对费米子场有
\begin{equation}
   \left\langle\chi(t,x)\Bar{\chi}(s,y)\right\rangle
   =\int_p\mathrm{e}^{ip(x-y)}\frac{\mathrm{e}^{-(t+s)p^2}}{i\Slash{p}+m_0}.
\label{eq:(3.18)}
\end{equation}
除了这些“量子传播子”之外，
式~\eqref{eq:(3.4)}、\eqref{eq:(3.12)} 和~\eqref{eq:(3.13)} 的微扰展开中还出现热核，
即式~\eqref{eq:(3.6)} 与~\eqref{eq:(3.14)}。

现在说明流动场微扰展开（流费曼图）的图形表示。对量子传播子%
~\eqref{eq:(3.17)} 与~\eqref{eq:(3.18)}，我们采用标准约定：规范玻色子的自由传播子用波浪线
表示，费米子的自由传播子用带箭头的直线表示。我们坚持这些约定，因为它们在包含费米子的系统
中十分自然。另一方面，文献~\cite{Suzuki:2013gza}和~\cite{Luscher:2011bx}
中曾用带箭头直线表示规范玻色子热核~\eqref{eq:(3.6)}；由于本文中该线型已被用于费米子传播
子，我们改用“双线”来表示热核~\eqref{eq:(3.6)} 与~\eqref{eq:(3.14)}。例如，
图~\ref{fig:1}--\ref{fig:5} 是对流动费米子场两点函数有贡献的单圈流费曼图。在这些图中，
双直线代表费米子热核~\eqref{eq:(3.14)}；箭头标记的是费米子数的流向，而不是流时间的方向。
类似地，图~\ref{fig:15} 中的双波浪线是规范玻色子热核~\eqref{eq:(3.6)}。在这种表示下我们
丢失了流时间方向的信息——该信息在文献~\cite{Suzuki:2013gza}
和~\cite{Luscher:2011bx} 中由箭头表示——不过这一信息很容易追溯回来。

流费曼图的另一要素是顶点。源自原始作用量~\eqref{eq:(2.1)} 的顶点用实心圆表示；经由流方
程、即式~\eqref{eq:(3.7)}、\eqref{eq:(3.15)} 与~\eqref{eq:(3.16)} 而来的顶点则如
图~\ref{fig:1}--\ref{fig:5} 中那样用空心圆表示；这些关于顶点的约定与文
献~\cite{Suzuki:2013gza}和~\cite{Luscher:2011bx}中的完全一致。

\subsection{加圈费米子场}
流动场的一个显著特点是 UV 有限性：当以重正化参数表达时，流动规范场 $B_\mu(t,x)$（$t$
严格为正）的任何关联函数都\emph{无需\/}相乘（波函数）
重正化即 UV 有限~\cite{Luscher:2011bx}。而且，即使取到等点极限，该有限性依然成立。因此，
$B_\mu(t,x)$（$t>0$）的任何局域乘积的任何关联函数，除参数重正化外无需进一步重正化便是 UV
有限的。换言之，尽管这些局域乘积通过流方程由裸规范场~$A_\mu(x)$ 的某个组合给出，它们却是
重正化了的有限量。这种 UV 有限性的一个基本原因在于流动规范场的传播子~\eqref{eq:(3.17)}
含有高斯衰减因子~$\sim \mathrm{e}^{-tp^2}$，它对~$t>0$ 实际上提供了紫外截断。但要证明上
述有限性，还必须利用该系统内在的、相对于规范势为\emph{非齐次\/}的
Becchi--Rouet--Stora（BRS）对称性~\cite{Luscher:2011bx}。

遗憾的是，上述第一层意义的有限性对流费米子场并不成立：它\emph{需要\/}波函数重正化。尽管其
传播子~\eqref{eq:(3.18)} 也带有衰减因子~$\sim \mathrm{e}^{-tp^2}$，但 BRS 变换作用在费米
子场（以及一般物质场）上是\emph{齐次\/}的，文献~\cite{Luscher:2011bx}中的有限性证明不适用
于它。事实上，对图~\ref{fig:1}--\ref{fig:5}（及箭头方向相反的图）的单圈图的计算表明，波函
数重正化
\begin{equation}
   \chi_R(t,x)=Z_\chi^{1/2}\chi(t,x),\qquad
   \Bar{\chi}_R(t,x)=Z_\chi^{1/2}\Bar{\chi}(t,x),\qquad
   Z_\chi=1+\frac{g^2}{(4\pi)^2}C_2(R)3\frac{1}{\epsilon}+O(g^4)
\label{eq:(3.19)}
\end{equation}
使关联函数成为 UV 有限~\cite{Luscher:2013cpa}。上面第二层意义上的有限性仍然成立：
$\chi_R(t,x)$ 与~$\Bar{\chi}_R(t,x)$ 的任何局域乘积的任何关联函数保持 UV
有限~\cite{Luscher:2013cpa}。

\begin{figure}
\begin{minipage}{0.3\textwidth}
\begin{center}
\includegraphics[width=3cm]{images/C01}
\caption{C01}
\label{fig:1}
\end{center}
\end{minipage}
\begin{minipage}{0.3\textwidth}
\begin{center}
\includegraphics[width=3cm]{images/C02}
\caption{C02}
\label{fig:2}
\end{center}
\end{minipage}
\begin{minipage}{0.3\textwidth}
\begin{center}
\includegraphics[width=3cm]{images/C03}
\caption{C03}
\label{fig:3}
\end{center}
\end{minipage}
\end{figure}

\begin{figure}
\begin{minipage}{0.3\textwidth}
\begin{center}
\includegraphics[width=3cm]{images/C04}
\caption{C04}
\label{fig:4}
\end{center}
\end{minipage}
\begin{minipage}{0.3\textwidth}
\begin{center}
\includegraphics[width=3cm]{images/C05}
\caption{C05}
\label{fig:5}
\end{center}
\end{minipage}
\end{figure}

尽管第二层意义上的有限性对我们的目的相当有用，我们仍须把重正化因子~$Z_\chi$（%
式~\eqref{eq:(3.19)}）纳入处理。这给问题带来了复杂性，因为我们必须在维数正规化与格点正规
化的 $Z_\chi$ 之间找到匹配因子。

避免这一复杂性的一个可能做法是用费米子动能算符的真空期望值来归一化费米子场：%
\footnote{在动能算符中，对 $N_{\mathrm{f}}$ 个费米子味道的求和已理解为完成。在本论文的第
一版中，我们曾用标量凝聚来归一化费米子场；该选择会带来与无质量费米子相联系的另一复杂性，
看来使用动能算符合适得多。}
\begin{align}
   \mathring{\chi}(t,x)
   &=\sqrt{\frac{-2\dim(R)N_{\mathrm{f}}}
   {(4\pi)^2t^2
   \left\langle\Bar{\chi}(t,x)\overleftrightarrow{\Slash{D}}\chi(t,x)
   \right\rangle}}
   \,\chi(t,x),
\label{eq:(3.20)}\\
   \mathring{\Bar{\chi}}(t,x)
   &=\sqrt{\frac{-2\dim(R)N_{\mathrm{f}}}
   {(4\pi)^2t^2
   \left\langle\Bar{\chi}(t,x)\overleftrightarrow{\Slash{D}}\chi(t,x)
   \right\rangle}}
   \,\Bar{\chi}(t,x),
\label{eq:(3.21)}
\end{align}
其中
\begin{equation}
   \overleftrightarrow{D}_\mu\equiv D_\mu-\overleftarrow{D}_\mu,
\label{eq:(3.22)}
\end{equation}
从而相乘重正化因子~$Z_\chi$ 在新的加圈变量中被消去。注意 $\mathring{\chi}(t,x)$
和~$\mathring{\Bar{\chi}}(t,x)$ 的质量量纲对任意~$D$ 都是 $3/2$，而 $\chi(t,x)$
和~$\Bar{\chi}(t,x)$ 的是~$(D-1)/2$。

动能算符真空期望值的最低阶（单圈）近似由图~\ref{fig:6} 中的图 D01 给出。对
$D=4-2\epsilon$ 维，我们有
\begin{equation}
   \left\langle\Bar{\chi}(t,x)\overleftrightarrow{\Slash{D}}\chi(t,x)
   \right\rangle
   =\frac{-2\dim(R)N_{\mathrm{f}}}{(4\pi)^2t^2}(8\pi t)^\epsilon\left[1+O(m_0^2t)\right].
\label{eq:(3.23)}
\end{equation}
注意在本设定中，真空期望值所需的质量标度由流时间~$t$ 提供。得到该表达式后，
式~\eqref{eq:(3.20)} 和~\eqref{eq:(3.21)} 中的常数因子已被选成：对足够小的流时间（%
$m_0^2t\ll1$），加圈变量与原始变量之差为~$O(g_0^2)$。

\begin{figure}
\begin{center}
\includegraphics[width=3cm]{images/D01}
\caption{D01}
\label{fig:6}
\end{center}
\end{figure}

\begin{figure}
\begin{minipage}{0.3\textwidth}
\begin{center}
\includegraphics[width=3cm]{images/D02}
\caption{D02}
\label{fig:7}
\end{center}
\end{minipage}
\begin{minipage}{0.3\textwidth}
\begin{center}
\includegraphics[width=3cm]{images/D03}
\caption{D03}
\label{fig:8}
\end{center}
\end{minipage}
\begin{minipage}{0.3\textwidth}
\begin{center}
\includegraphics[width=3cm]{images/D04}
\caption{D04}
\label{fig:9}
\end{center}
\end{minipage}
\end{figure}

\begin{figure}
\begin{minipage}{0.3\textwidth}
\begin{center}
\includegraphics[width=3cm]{images/D05}
\caption{D05}
\label{fig:10}
\end{center}
\end{minipage}
\begin{minipage}{0.3\textwidth}
\begin{center}
\includegraphics[width=3cm]{images/D06}
\caption{D06}
\label{fig:11}
\end{center}
\end{minipage}
\begin{minipage}{0.3\textwidth}
\begin{center}
\includegraphics[width=3cm]{images/D07}
\caption{D07}
\label{fig:12}
\end{center}
\end{minipage}
\end{figure}

\begin{figure}
\begin{center}
\includegraphics[width=3cm]{images/D08}
\caption{D08}
\label{fig:13}
\end{center}
\end{figure}

\begin{table}
\caption{各图对式~\eqref{eq:(3.24)} 的贡献，单位为
$\frac{-2\dim(R)N_{\mathrm{f}}}{(4\pi)^2t^2}\frac{g_0^2}{(4\pi)^2}C_2(R)$。}
\label{table:1}
\begin{center}
\renewcommand{\arraystretch}{2.2}
\setlength{\tabcolsep}{20pt}
\begin{tabular}{cr}
\toprule
 diagram & \\
\midrule
 D02 & $-\dfrac{1}{\epsilon}-2\ln(8\pi t)+O(m_0^2t)$ \\
 D03 & $2\dfrac{1}{\epsilon}+4\ln(8\pi t)+2+4\ln2-2\ln3+O(m_0^2t)$ \\
 D04 & $-20\ln2+16\ln3+O(m_0^2t)$ \\
 D05 & $12\ln2-5\ln3+O(m_0^2t)$ \\
 D06 & $-4\dfrac{1}{\epsilon}-8\ln(8\pi t)-2+O(m_0^2t)$ \\
 D07 & $8\ln2-4\ln3+O(m_0^2t)$ \\
 D08 & $-2\ln3+O(m_0^2t)$ \\
\bottomrule
\end{tabular}
\end{center}
\end{table}

期望值~\eqref{eq:(3.23)} 的次领头阶（即双圈）表达式由图~\ref{fig:7}--\ref{fig:13}
（及箭头方向相反的图）中的流费曼图给出。这些图的计算稍嫌繁杂，但可以用与文
献~\cite{Luscher:2010iy}附录 B 类似的方式完成［配合我们附录~\ref{sec:B} 的积分公式，
式~\eqref{eq:(B1)} 与~\eqref{eq:(B2)}］。每个图的贡献列于表~\ref{table:1}。合计得
\begin{align}
   &\left\langle\Bar{\chi}(t,x)\overleftrightarrow{\Slash{D}}\chi(t,x)
   \right\rangle
\notag\\
   &=\frac{-2\dim(R)N_{\mathrm{f}}}{(4\pi)^2t^2}
   \left\{(8\pi t)^\epsilon
   +\frac{g_0^2}{(4\pi)^2}C_2(R)
   \left[-3\frac{1}{\epsilon}-6\ln(8\pi t)+\ln(432)\right]
   +O(m_0^2t)+O(g_0^4)\right\}.
\label{eq:(3.24)}
\end{align}
利用式~\eqref{eq:(A1)} 处理式~\eqref{eq:(3.20)} 和~\eqref{eq:(3.21)}
中的归一化因子，可得
\begin{equation}
   \frac{-2\dim(R)N_{\mathrm{f}}}
   {(4\pi)^2t^2\left\langle\Bar{\chi}(t,x)
   \overleftrightarrow{\Slash{D}}\chi(t,x)\right\rangle}
   =Z(\epsilon)
   \left\{1+\frac{g^2}{(4\pi)^2}C_2(R)\left[3\frac{1}{\epsilon}-\Phi(t)\right]
   +O(m^2t)+O(g^4)\right\},
\label{eq:(3.25)}
\end{equation}
其中
\begin{equation}
   Z(\epsilon)\equiv\frac{1}{(8\pi t)^\epsilon},
\label{eq:(3.26)}
\end{equation}
以及
\begin{equation}
   \Phi(t)\equiv-3\ln(8\pi\mu^2t)+\ln(432).
\label{eq:(3.27)}
\end{equation}
